% =============================================================================
% Section III — Dataset and Preprocessing
% =============================================================================
\section{Dataset and Preprocessing}
\label{sec:dataset}

\subsection{Raw Data Overview}
The project uses a multivariate water-network dataset with 6{,}000
time-stamped records collected from ten spatially distributed sensors
(S001--S010) covering a three-week period from 2024-01-01 to 2024-01-21
at an approximate 5-minute sampling interval. Each record includes:
sensor ID, UTC timestamp, pressure~(bar), flow rate~(L/s),
temperature~(\textdegree C), binary leak status, and binary burst status.

Burst events are excluded from the modelling corpus because they represent
a qualitatively different and more severe hydraulic anomaly. After burst
removal, 5{,}590 records remain with 774 labelled leak timesteps
(13.8\,\%), confirming moderate class imbalance at the raw record level
that becomes more pronounced at the window level after temporal splitting.

\subsection{Feature Exploration}
Figure~\ref{fig:pressure_dist} and Figure~\ref{fig:flow_dist} show
stratified distributions for pressure and flow rate. Leak records present
systematically lower pressure (shifted distribution towards 1.5--2.5\,bar)
while normal records concentrate around 3.0--3.5\,bar. Flow rate shows a
complementary rightward shift for leak windows, consistent with the
hydraulic inverse relationship between pressure loss and discharge increase.

\begin{figure}[H]
\centering
\includegraphics[width=\columnwidth]{images/fig1_pressure_distribution.jpg}
\caption{Pressure distribution for normal and leak classes. Leak records
exhibit a clear leftward shift towards lower pressures.}
\label{fig:pressure_dist}
\end{figure}

\begin{figure}[H]
\centering
\includegraphics[width=\columnwidth]{images/fig2_flowrate_distribution.jpg}
\caption{Flow-rate distribution for normal and leak classes. Leak records
show elevated flow relative to normal operating conditions.}
\label{fig:flow_dist}
\end{figure}

Temporal context is essential. Figure~\ref{fig:timeseries} illustrates that
leak signals manifest as sequential patterns spanning multiple consecutive
readings, not as isolated outlier points.

\begin{figure}[H]
\centering
\includegraphics[width=\columnwidth]{images/fig3_timeseries_leak_sensor.jpg}
\caption{Sensor S003 time series with confirmed leak events marked in red.
Pressure dips and flow spikes co-occur, reinforcing the need for
multivariate sequence modelling.}
\label{fig:timeseries}
\end{figure}

Correlation analysis (Figure~\ref{fig:correlation}) confirms that pressure
is the most informative raw feature, with a Pearson correlation of
$r = -0.606$ with the leak label. Flow rate shows a moderate positive
correlation ($r = 0.197$), while temperature is essentially uncorrelated
($r = 0.023$) and retained only as auxiliary context.

\begin{figure}[H]
\centering
\includegraphics[width=\columnwidth]{images/fig4_correlation_matrix.jpg}
\caption{Pearson correlation matrix. Pressure exhibits the strongest
association with leak status ($r = -0.606$).}
\label{fig:correlation}
\end{figure}

\subsection{Feature Engineering and Windowing}
Three domain-informed features are computed per record to supplement raw
sensor readings:
\begin{itemize}[leftmargin=*]
\item \textbf{Pressure delta} ($\Delta P_t = P_t - P_{t-1}$, per sensor):
captures rapid pressure drops at leak onset.
\item \textbf{Flow z-score} (sensor-normalised): removes sensor-specific
baseline offsets to highlight relative flow excursions.
\item \textbf{Pressure-to-flow ratio} ($r_t = P_t / (Q_t + \varepsilon)$):
collapses the inverse hydraulic relationship into a single discriminative scalar.
\end{itemize}

The final feature vector per timestep is therefore $\mathbf{x}_t \in
\mathbb{R}^{6}$, comprising pressure, flow rate, temperature, pressure delta,
flow z-score, and pressure-to-flow ratio.

Sliding windows are constructed independently per sensor with:
\begin{itemize}[leftmargin=*]
\item Window size: $W = 20$ timesteps (approximately 100\,minutes).
\item Step size: $S = 5$ timesteps (25\,minutes).
\item Label: positive if \emph{any} timestep in the window is a leak.
\end{itemize}

The temporal 80/20 split per sensor yields the following window counts:
\begin{itemize}[leftmargin=*]
\item \textbf{Training total}: 863 windows --- leak: 329 (38.1\,\%).
\item \textbf{Test set}: 221 windows --- leak: 91 (41.2\,\%).
\end{itemize}

From the 863 training windows, a stratified internal split produces a
\emph{fit set} (656 windows, 34.9\,\% leak, after ratio capping) and a
\emph{calibration set} (173 windows, 38.2\,\% leak) used exclusively
for early stopping and threshold selection. Feature standardisation
statistics are computed on the fit set only and applied identically to
calibration and test data.

